\documentclass{article}

\usepackage[T1]{fontenc}
\usepackage{imakeidx}
\usepackage[pdfborder={0 0 0},bookmarks=true,bookmarksnumbered=true,breaklinks=true]{hyperref}
\usepackage{bookmark}
\usepackage{longtable}
\usepackage{upquote}
\DeclareUnicodeCharacter{2194}{\ensuremath{\leftrightarrow}}

% Include auto-generated fragments from bch-docgen.
% Run: cargo run -p bch-docgen
\newcommand{\bchdoc}[1]{\input{doc/generated/#1}}
\providecommand{\bchdocversion}{unknown}
\InputIfFileExists{doc/generated/build-version.tex}{}{}

\title{bcachefs: Principles of Operation}
\author{Kent Overstreet (and others)}

\date{Revision \texttt{\bchdocversion}\\\today}

\begin{document}

\maketitle
\setcounter{tocdepth}{4}
\setcounter{secnumdepth}{4}
\tableofcontents

\section{Introduction and overview}

Bcachefs is a modern, general purpose, copy on write filesystem descended from
bcache, a block layer cache.

The internal architecture is very different from most existing filesystems where
the inode is central and many data structures hang off of the inode. Instead,
bcachefs is architected more like a filesystem on top of a relational database,
with tables for the different filesystem data types - extents, inodes, dirents,
xattrs, et cetera.

The entire filesystem is built on two primitives: \hyperref[sec:btrees]{btrees} and buckets. All
filesystem state is a key-value pair in a btree. A handful of btrees store core
filesystem data (extents, inodes, dirents, xattrs, subvolumes), while others
handle allocation tracking, reverse mappings, background maintenance, and crash
recovery. There are no separate inode tables, bitmap allocators, or per-inode
extent trees. On the storage side, all disk space is divided into buckets with
generation numbers, described in the allocation section below. The btrees map
logical filesystem objects to physical bucket locations, and backpointers
provide the reverse mapping. Everything else in bcachefs is built on this
transactional key-value store over generational bucket storage.

bcachefs supports almost all of the same features as other modern COW
filesystems, such as ZFS and btrfs, but in general with a cleaner, simpler,
higher performance design.

\subsection{Performance overview}

The core of the architecture is a very high performance and very low latency b+
tree, which also is not a conventional b+ tree but more of hybrid, taking
concepts from compacting data structures: btree nodes are very large, log
structured, and compacted (resorted) as necessary in memory. This means our b+
trees are very shallow compared to other filesystems. Other COW filesystems use
the Linux page cache for metadata, which limits them to 4K btree nodes; bcachefs
manages its own btree node cache, so we can use 128K--256K nodes and tune
reclaim independently of the page cache. At petabyte scale with spinning disks,
deep btrees with 4K nodes mean more seeks per lookup and level-2 nodes too large
to stay resident under page cache pressure.

What this means for the end user is that since we require very few seeks or disk
reads, filesystem latency is extremely good - especially cache cold filesystem
latency, which does not show up in most benchmarks but has a huge impact on real
world performance, as well as how fast the system "feels" in normal interactive
usage. Latency has been a major focus throughout the codebase - notably, we have
assertions that we never hold b+ tree locks while doing IO, and the btree
transaction layer makes it easy to aggressively drop and retake locks as
needed - one major goal of bcachefs is to be the first general purpose soft
realtime filesystem.

Additionally, unlike other COW btrees, btree updates are journalled. This
greatly improves our write efficiency on random update workloads, as it means
btree writes are only done when we have a large block of updates, or when
required by memory reclaim or journal reclaim.

\subsection{Transactions}

All btree operations go through a transaction layer (\texttt{btree\_trans}) that
provides atomic multi-key updates with optimistic concurrency control.
Transactions collect updates in memory, then commit them atomically: first
acquiring a journal reservation, then taking write locks on affected btree nodes
in sorted order, applying updates, and releasing locks after the journal has
accepted the change (but before the journal write is persisted to disk). This
means btree locks are held only for in-memory operations, never across IO.

Btree nodes use six-locks (shared/intent/exclusive): the intent state allows an
operation that modifies multiple nodes to hold intent locks for its duration,
upgrading to write only for each individual in-memory update, avoiding the need
to hold write locks across entire split or merge operations.

Deadlock avoidance uses standard database cycle detection: before sleeping on a
contended lock, the transaction checks for lock cycles among all waiting
transactions. If a cycle is found, one transaction drops all its locks and
restarts from the beginning. This avoids deadlocks entirely and keeps
worst-case latency bounded, at the cost of
requiring all transaction code to be idempotent and prepared for restarts at any
point. This idempotency requirement has a profound consequence for reliability:
since every operation can be safely interrupted and restarted, the filesystem is
inherently resilient to interruption at any point - including during recovery
itself. A crash during journal replay or an on-disk format upgrade simply
results in the operation continuing from a safe point on the next mount.
Recovery passes generally restart from the beginning and re-check everything;
only passes that maintain persistent work queues can resume mid-pass.

Sequence numbers on btree nodes enable efficient relocking after a restart: if a
node's sequence number hasn't changed, the lock can be reacquired without
re-traversing the btree.

Transactions also support triggers: callbacks that fire on key insertion or
deletion. Transactional triggers run before commit and may generate additional
updates (e.g. updating backpointers when an extent changes). Atomic triggers
run with the journal lock held and are used for accounting updates that must be
exactly synchronized with the commit.

For operations that cannot be made atomic within a single transaction (such as
truncate or file insert/collapse, which may need to modify an unbounded number
of extents), bcachefs uses logged operations: the intent is written to the
\texttt{logged\_ops} btree before starting, and on recovery any incomplete
logged operations are replayed to completion.

\subsection{Recovery}

On unclean shutdown, the \hyperref[sec:journal]{journal} is replayed to bring the btrees up to date.
Journal entries record btree updates in commit order; replay applies them
sequentially to reconstruct the state at the last completed commit.

To guarantee ordering of btree updates after a crash, we need to detect when a
btree node was flushed to disk but its corresponding journal entry was not;
otherwise, recovery could see later updates without the earlier updates they
depended on.

During btree node reads, bcachefs detects this situation: bsets (sorted key
arrays within a node) whose journal sequence number is newer than the last
successfully written journal entry indicate the btree node was written to disk
but the corresponding journal entry was not. These sequences are blacklisted
in the \hyperref[sec:superblock]{superblock} (\texttt{bch\_sb\_field\_journal\_seq\_blacklist}), and
bsets referencing them are ignored until the btree node is next rewritten.
After an unclean shutdown, a padding of 64 sequence numbers past the last
journal entry read is also blacklisted, to cover in-flight btree writes.
Blacklisted sequences are never reused---the 64-bit counter simply skips
past them. Blacklist entries are garbage collected once no btree node on disk
references them.

\subsubsection{Recovery passes}

Beyond journal replay, bcachefs has approximately 48 recovery passes with an
explicit dependency graph. Not all passes run on every mount - which passes run
depends on context:

\begin{description}
	\item[Every mount] A default set of passes is scheduled: reading
		accounting, bucket generations, and snapshots first, checking
		snapshot consistency, then resuming logged operations and
		cleaning up dead inodes.

	\item[Unclean shutdown] Adds journal replay. No extra repair passes
		are needed because the filesystem is always consistent after
		journal replay (absent bugs or data corruption).

	\item[Version upgrade/downgrade] Each version transition defines
		specific passes to run for format migration.

	\item[Fsck] A defined subset of validation passes runs: checking inodes,
		extents, dirents, xattrs, link counts, and directory structure.
		Note that \texttt{check\_backpointers\_to\_extents} (the reverse
		direction check) is too expensive to run during normal fsck.
\end{description}

Recovery is split into two phases around the transition to read-write mode.
Early passes run in-memory only, accumulating delayed metadata writes: reading
btrees into memory, verifying btree topology, and reconstructing allocation
state. Once the allocator and journal are initialized, the filesystem
transitions to read-write and journal replay runs - followed by the heavier
validation and repair passes that need to write.

The filesystem can schedule new recovery passes at any time, including during
other recovery passes or normal operation, as soon as it notices something needs
repair. In early recovery, this can rewind the recovery sequence to run earlier
passes again. Most validation passes (those marked \texttt{PASS\_ONLINE}) can
also run after the filesystem is fully mounted, either triggered explicitly or
scheduled automatically when an inconsistency is detected at runtime.

For catastrophic damage, heavier passes exist that reconstruct entire btrees
from scratch. These are triggered automatically when normal recovery detects
structural damage: \texttt{scan\_for\_btree\_nodes} is scheduled when
\texttt{check\_topology} finds missing or unreachable btree nodes - it
physically scans the raw device for valid btree node headers and recovers what
it can. \texttt{reconstruct\_snapshots} is scheduled when the snapshots btree
itself is lost, or when references to missing snapshots cannot be resolved
by deletion alone and recovering those snapshots is feasible.
\texttt{check\_allocations} rebuilds all allocation accounting from the extents
and backpointers btrees. These passes are not part of normal recovery and may
take considerable time on large filesystems.

When a pass is scheduled at runtime (post-mount), its persistence depends on
severity. Passes scheduled without the \texttt{nopersistent} flag are written
to the superblock's \texttt{recovery\_passes\_required} field and survive
reboots: the pass will run automatically at the next mount even if the
filesystem is cleanly unmounted. Ephemeral scheduling (used when a prior pass
must complete first before the repair makes sense) is remembered only until
the current recovery sequence completes and does not touch the superblock.
Once a pass runs successfully, it clears its own bit from
\texttt{recovery\_passes\_required}.

Some passes are scheduled with a ratelimit flag, used for errors that may be
detected in large numbers during a single I/O operation (for example, EC
accounting inconsistencies during a stripe write). A ratelimited pass is not
re-scheduled unless the same error is detected again via a non-ratelimited
request, preventing repair-loop flooding. The ratelimiting state is in-memory
only and does not affect persistence.

The complete list of recovery passes, in execution order:

\bchdoc{recovery-passes}

\subsection{Durability and fsync semantics}

Because all metadata updates pass through a single, totally ordered journal,
bcachefs provides a particularly simple crash recovery contract: \emph{prefix
consistency}. The state recovered after a crash is exactly the state of the
filesystem as of some recent moment $T$: every operation that completed before
$T$ is present, every operation after $T$ is absent, and no other outcomes are
possible. There is no reordering and no partial visibility --- recovery never
shows operation $B$ without an operation $A$ that preceded it, whether or not
the two were related. This is strictly stronger than the ``strictly ordered
metadata consistency'' (SOMC) provided by ext4 and XFS, which promises ordering
only between \emph{dependent} operations.

\texttt{fsync(2)} constrains $T$: it completes the file's outstanding data
writeback, then flushes the journal through the inode's most recent
modification, guaranteeing $T$ is at least that point. Together with prefix
consistency this gives fsync the following concrete meaning:

\begin{itemize}
\item Everything that ever modified the fsynced inode is persisted: writes,
truncates, timestamps, extended attributes --- and link count changes, so an
fsync of a file persists all of that file's hard link directory entries
(creating or removing a hard link updates the inode in the same transaction as
the directory entry).
\item All operations system-wide that preceded the inode's last modification
are also persisted, related to the fsynced file or not.
\item An fsync of a directory persists all entry creations, removals and
renames in that directory up to its last modification (directory entry
updates modify the directory inode in the same transaction).
\end{itemize}

What is \emph{not} guaranteed: operations that occurred \emph{after} the
fsynced inode's last modification, even if they preceded the fsync call
itself. If file A is modified, then file B is modified, then A is fsynced,
B's modification is not covered --- it will often be persisted anyway, because
journal writes carry everything up to the journal's current position, but this
is an accident of batching, not a promise. (Some filesystems' fsync
implementations track fine-grained dependencies and may persist more, or less,
in this situation; a number of xfstests in the crash-consistency family encode
one such implementation's behavior rather than any standardized contract.)

Data writes participate in the prefix property in normal (copy on write)
operation: an extent is never journalled until the write of its data has
completed, so if recovery shows an extent, the data it points to is intact.

Exceptions, both opt-in:

\begin{itemize}
\item \texttt{lazytime}: timestamp-only updates are deliberately buffered in
memory and may be lost on crash, up to the lazytime persistence points (sync,
fsync, inode eviction, periodic expiry).
\item \texttt{nocow}: data is overwritten in place, so the data prefix
property is weakened --- after a crash, a nocow extent may contain newer data
than moment $T$, and without checksums this cannot be detected.
\end{itemize}

The common write-temporary-then-rename-over pattern gets a stronger
guarantee: when a rename \emph{overwrites} an existing file, bcachefs first
completes writeback of the source file's data, so the data's extents are
journalled before the rename is. Combined with prefix consistency this makes
the replacement atomic and complete without any fsync: after a crash, either
the old file is intact, or the rename happened and all of the new file's data
is present --- the empty-or-truncated-file outcome cannot occur. Note this
orders the rename behind the data but does not make the rename itself durable
(no journal flush is issued); an application that must know the replacement
has happened still needs fsync. Plain renames that don't overwrite anything
get no implicit writeback.

Related interfaces, briefly: \texttt{fdatasync(2)} is implemented
identically to \texttt{fsync(2)} (stronger than POSIX requires).
\texttt{syncfs(2)} and \texttt{sync(2)} flush the entire journal, making
every operation completed before the call durable. \texttt{O\_DIRECT}
confers no durability: completion means the data has been transferred to the
device, but the extent describing it is journalled metadata like any other,
so until a journal flush a crash simply recovers to a moment before the
write. Crash semantics are identical to a buffered write whose writeback has
completed; fsync or fdatasync is still what makes the write survive. The \texttt{O\_TMPFILE}-then-\texttt{linkat} pattern is
covered by fsync's own-inode completeness: the link updates the inode, so a
subsequent fsync of the file persists it. Finally, truncate-then-rewrite of an existing file
(\texttt{O\_TRUNC}) is inherently unsafe as a replacement strategy on any
filesystem: the old contents are destroyed before the new contents exist, so
no placement of fsync can provide an old-or-new guarantee --- a crash after
the truncate persists but before the rewrite does recovers an empty file,
and fsync only narrows that window. Atomic replacement is what
rename-over is for.

\subsection{Bucket based allocation}

As mentioned, bcachefs is descended from bcache, where the ability to efficiently
invalidate cached data and reuse disk space was a core design requirement. To
make this possible the allocator divides the disk up into buckets, typically
512k to 2M but possibly larger or smaller. Buckets and data pointers have
generation numbers: we can reuse a bucket with cached data in it without finding
and deleting all the data pointers by incrementing the generation number. The
generation number mechanism was originally created for invalidating cached data,
but it turned out to be valuable well beyond that: it makes bootstrapping
allocation information during journal replay much more tractable, and it makes
it possible to repair certain kinds of corruption that would be unrecoverable
otherwise.

In keeping with the copy-on-write theme of avoiding update in place wherever
possible, we never rewrite or overwrite data within a bucket - when we allocate
a bucket, we write to it sequentially and then we don't write to it again until
the bucket has been invalidated and the generation number incremented.

This means we require a copying garbage collector to deal with internal
fragmentation, when patterns of random writes leave us with many buckets that
are partially empty (because the data they contained was overwritten) - copy GC
evacuates buckets that are mostly empty by writing the data they contain to new
buckets. This also means that we need to reserve space on the device for the
copy GC reserve when formatting - 8\% by default (configurable from 5\% to
21\%). A fragmentation LRU btree tracks bucket fill levels so that copygc can
find the most fragmented buckets directly instead of scanning the entire alloc
btree.

The allocator originally used in-memory bucket arrays and dedicated scanning
threads to maintain free lists, discard lists, and eviction heaps. These were
replaced by btree-based data structures: a freespace btree, a discard btree,
and an LRU btree. The scanning threads are gone entirely; the code that
replaces them is transactional btree code, much of it trigger-based, which is
far easier to debug and reason about. The old threads were also prone to
excessive CPU usage when the filesystem was nearly full; the btree approach
eliminated those corner cases.

There are some advantages to structuring the allocator this way, besides being
able to support cached data:
\begin{itemize}
	\item By maintaining multiple \hyperref[sec:write-points]{write points} that are writing to different buckets,
		we're able to easily and naturally segregate unrelated IO from different
		processes, which helps greatly with fragmentation.

	\item The fast path of the allocator is essentially a simple bump allocator - the
		disk space allocation is extremely fast

	\item Fragmentation is generally a non issue unless copygc has to kick
		in, and it usually doesn't under typical usage patterns. The
		allocator and copygc are doing essentially the same things as
		the flash translation layer in SSDs, but within the filesystem
		we have much greater visibility into where writes are coming
		from and how to segregate them, as well as which data is
		actually live - performance is generally more predictable than
		with SSDs under similar usage patterns.

	\item The same algorithms are intended to be used for managing SMR
		(shingled magnetic recording) and zoned hard drives directly,
		avoiding the translation layer in the drive. Buckets map
		naturally to zones: when we allocate a bucket, we write to it
		once in an append-only fashion, then never write to it again
		until we discard and reuse the whole thing --- exactly the
		semantics a zoned device requires.
\end{itemize}

\subsection{Backpointers}

Every sector range on disk that contains data or metadata has a corresponding
backpointer: a reverse reference from the physical location back to the logical
btree entry that owns it. Backpointers enable efficient answers to the question
``what data lives in this bucket?'' without scanning the entire extents btree.
This is essential for copygc, device evacuation, scrub, and the reconcile
subsystem.

See \hyperref[sec:backpointers]{Backpointers} in subsystem details for
on-disk structure, maintenance, and consistency validation.

\section{Feature overview}

\subsection{IO path options}
\label{sec:io-path-options}

Most options that control the IO path can be set at either the filesystem level
or on individual inodes (files and directories). When set on a directory via the
\texttt{bcachefs set-file-option} command, they will be automatically applied
recursively.

Checksumming, compression, and encryption all operate at the granularity of
entire extents (up to \texttt{encoded\_extent\_max}, default 256 KB). An extent
that has been through any of these transformations is called an \emph{encoded
extent}. Because the checksum covers the full encoded extent, any read within it
--- even a single sector --- must read, verify, and (if compressed) decompress
the whole thing. This is the fundamental performance tradeoff shared by all
three features: they improve data integrity and storage efficiency at the cost
of read amplification for small random IO. For workloads dominated by small
random reads (databases, VMs), reducing \texttt{encoded\_extent\_max} or
disabling these features may be appropriate.

\subsubsection{Checksumming}

bcachefs supports both metadata and data checksumming --- crc32c by default, but
stronger checksums are available as well. Enabling data checksumming incurs some
performance overhead: besides the checksum calculation, writes have to be
bounced for checksum stability (Linux generally cannot guarantee that the buffer
being written is not modified in flight), but reads generally do not have to be
bounced.

\subsubsection{Encryption}

bcachefs supports authenticated (AEAD) encryption using ChaCha20/Poly1305.
Unlike typical block layer or filesystem encryption (usually AES-XTS), which
operates on fixed blocks with no room for nonces or MACs, bcachefs stores a
nonce and cryptographic MAC alongside every data pointer. This creates a chain
of trust from the \hyperref[sec:superblock]{superblock} (or \hyperref[sec:journal]{journal}) down to individual extents: any
modification, deletion, reordering, or rollback of metadata is detectable.

Encryption is all-or-nothing at the filesystem level: all data and metadata
except the superblock is encrypted, and all data and metadata is authenticated.
Per-file or per-directory encryption is not supported because btree nodes are
shared structures that do not belong to any individual file. Encryption can
only be enabled at format time; it cannot be added to an existing filesystem.

\textbf{Critical exception}: Data written with the \texttt{nocow} option is
stored \textbf{unencrypted}, even on an encrypted filesystem. This is a hard
design incompatibility, not a policy choice: ChaCha20 requires a unique nonce
per (key, plaintext), and bcachefs stores the nonce externally alongside each
data pointer. An in-place \texttt{nocow} overwrite cannot produce a new nonce
without also atomically updating the pointer (which defeats the purpose of
nocow), so the (key, nonce) pair would be reused against a different
plaintext --- catastrophic for ChaCha20, since the XOR of two ciphertexts
leaks the XOR of the plaintexts. Supporting encrypted nocow would require
integrating a position-tweakable cipher (e.g.\ AES-XTS) alongside ChaCha20.
See the Nocow section below and the Nonce reuse discussion for related
considerations.

The key hierarchy uses a passphrase (never stored on disk) to derive a
passphrase key via scrypt, which decrypts a random master key stored in the
superblock. The master key encrypts all data and metadata. Changing the
passphrase re-encrypts only the master key, not filesystem data. The KDF runs
in userspace, so alternative key sources (hardware tokens, key files) can be
integrated without kernel changes. There is currently no key rotation, key
escrow, or multi-passphrase support. Losing the passphrase means permanent
data loss.

At mount time, the kernel retrieves the passphrase key from the Linux kernel
keyring (\texttt{bcachefs:<UUID>}). Note that keys added to a session keyring
are not visible from other sessions --- use \texttt{KEY\_SPEC\_USER\_KEYRING}
for cross-session visibility. Similarly, \texttt{sudo mount} uses root's
keyring, not the calling user's.

By default, MACs are truncated to 80 bits, which is sufficient for most threat
models. The \texttt{wide\_macs} option stores the full 128-bit MAC and is
recommended when the storage device itself is untrusted (e.g. USB drives,
network storage). Metadata always uses full-width MACs.

\paragraph{Nonce reuse}

ChaCha20 requires a unique nonce for every (key, plaintext) pair. bcachefs
stores nonces externally (in the data pointer), so the filesystem itself
assigns nonces --- and anything that breaks the ``one nonce per plaintext''
invariant breaks encryption security. There are two distinct ways this can
happen:

\emph{External snapshots, mounted read-write}: if the underlying storage is
snapshotted externally (LVM, ZFS zvol, VM snapshot) and the snapshot is
mounted read-write, both instances share the same master key and will assign
nonces independently, producing collisions. bcachefs's own snapshots do not
have this problem.

\textbf{Mitigation}: Never mount an external snapshot of an encrypted volume
read-write --- keep external snapshots read-only (\texttt{-o nochanges}).
Alternatively, place LUKS between the snapshot layer and bcachefs (e.g.
LVM $\to$ LUKS $\to$ bcachefs).

\emph{In-place overwrites (\texttt{nocow})}: in-place writes cannot update
the externally-stored nonce without also updating the data pointer, which
would defeat the point of nocow. Encrypting a nocow overwrite would therefore
reuse the existing (key, nonce) against new plaintext. This is why nocow
data on encrypted filesystems is stored in plaintext; see the Nocow section
for details.

See \hyperref[sec:encryption]{Encryption} in subsystem details for key
hierarchy internals, keyring integration, and MAC storage details.

\subsubsection{Compression}

bcachefs supports gzip, lz4 and zstd compression, operating at extent
granularity as described above. This gives better compression ratios than
filesystems that compress at the 4K page level: compression algorithms find more
redundancy when fed more data at once, and rounding compressed output back up to
block alignment at small granularity loses much of the compression ratio.
Per-extent granularity also simplifies the IO path: everything works in terms of
extents, and there is nothing smaller than an extent to write code to handle.

Data can also be compressed or recompressed with a different algorithm in the
background by the reconcile subsystem, if the \texttt{background\_compression}
option is set.

\subsubsection{Nocow}

The \texttt{nocow} option enables in-place writes: data is overwritten directly
rather than being written to a new location (copy-on-write). This bypasses the
normal encoded-extent pipeline --- nocow data is not checksummed, not compressed,
and not encrypted, even on an encrypted filesystem.

Nocow mode is useful for workloads that require stable disk offsets or that
cannot tolerate the write amplification of COW (databases, VM images). It can
be set per-file or per-directory.

\hyperref[sec:snapshots]{Snapshots} and \hyperref[sec:reflink]{reflink} still trigger COW semantics even for nocow files --- when
a nocow file has shared extents (via snapshot or reflink), writes to those
extents will COW as usual. Once the shared reference is gone, subsequent writes
resume in-place.

Because nocow data has no checksums, it cannot be verified by scrub, cannot be
self-healed from a replica on checksum mismatch, and is not protected against
bitrot.

On encrypted filesystems, nocow data is stored in plaintext --- this is a
hard incompatibility, not an oversight. ChaCha20 requires unique nonces, and
bcachefs's nonces live in the data pointer (external to the ciphertext). An
in-place overwrite therefore cannot change the nonce, so it would reuse the
(key, nonce) pair against new plaintext, which leaks both plaintexts via
keystream recovery. This also means subsystems that re-encrypt or re-emit
data in the background (e.g.\ reconcile, background compression) must skip
nocow extents; treating a nocow extent as encrypted is a correctness bug.
Lifting this restriction would require adding a position-tweakable cipher
(e.g.\ AES-XTS) to the encryption pipeline.

These tradeoffs should be weighed carefully; for most workloads, the default
COW behavior with checksumming is the better choice.

\subsection{Multiple devices}

bcachefs is a multi-device filesystem. Devices need not be the same size: by
default, the allocator will \hyperref[sec:write-points]{stripe across all available devices} but biasing in
favor of the devices with more free space, so that all devices in the filesystem
fill up at the same rate. Devices need not have the same performance
characteristics: we track device IO latency and direct reads to the device that
is currently fastest.

\subsubsection{Replication}

bcachefs supports standard RAID1/10 style redundancy with the
\texttt{data\_replicas} and \texttt{metadata\_replicas} options. Layout is not
fixed as with RAID10: a given extent can be replicated across any set of
devices; the \texttt{bcachefs fs usage} command shows how data is replicated
within a filesystem.

Replica placement is flexible: each write picks devices from those available
for the requested data type, respecting target group constraints and balancing
across devices by free space. When a read detects a checksum error on one
replica, it transparently falls back to another replica and repairs the
corrupted copy by rewriting it. If the desired replica count is reduced (or
increased), the reconcile subsystem adjusts existing data in the background.

\subsubsection{Erasure coding}

bcachefs supports Reed-Solomon erasure coding, the same algorithm used by
most RAID5/6 implementations. Erasure coding is a per-inode option
(\texttt{erasure\_code}), so it can be enabled on specific directories via
\texttt{set-file-option} without affecting the rest of the filesystem. The
desired redundancy is taken from the \texttt{data\_replicas} option:
\texttt{data\_replicas=2} yields one parity block (RAID-5 style),
\texttt{data\_replicas=3} yields two (RAID-6 style), and
\texttt{data\_replicas=1} disables erasure coding entirely. Parity is limited
to at most two blocks. Erasure coding of metadata is not supported.

In conventional RAID, the ``write hole'' is a significant problem: a small write
within a stripe requires updating the P and Q (parity) blocks, and since those
writes cannot be done atomically, a crash can leave parity inconsistent,
corrupting reconstruct reads for unrelated data in the same stripe. ZFS avoids
this by making every write a new stripe, but the resulting fragmentation hurts
performance: reads of fragmented data are bounded by the latency of the slowest
fragment, driving median latency toward tail latency.

bcachefs takes advantage of copy-on-write to avoid both problems. Since
updating stripes in place is the root cause, we simply never do it. And since
per-extent stripes would cause the same fragmentation as ZFS, we erasure code
entire buckets, taking advantage of bucket-based allocation and copying
garbage collection.

See \hyperref[sec:erasure-coding]{Erasure coding} in subsystem details for
write path mechanics, stripe layout, fragmentation handling, and on-disk
representation.

\subsubsection{Device labels and targets}

By default, writes are striped across all devices in a filesystem, but they may
be directed to a specific device or set of devices with the various target
options. The allocator only prefers to allocate from devices matching the
specified target; if those devices are full, it will fall back to allocating
from any device in the filesystem.

Target options may refer to a device directly, e.g.
\texttt{foreground\_target=/dev/sda1}, or they may refer to a device label. A
device label is a path delimited by periods - e.g. ssd.ssd1 (and labels need not
be unique). This gives us ways of referring to multiple devices in target
options: If we specify ssd in a target option, that will refer to all devices
with the label ssd or labels that start with ssd. (e.g. ssd.ssd1, ssd.ssd2).

Device labels are used only for targeting. How replicas are spread across
devices to survive hardware failures is a separate, per-device property; see
\hyperref[sec:failure-domains]{Failure domains}.

Four target options exist. All may be set at the filesystem level (at format
time, at mount time, or at runtime via sysfs). All except
\texttt{metadata\_target} may also be set on a particular file or directory:

\begin{description}
	\item \texttt{foreground\_target}: normal foreground data writes, and
		metadata if \\ \texttt{metadata\_target} is not set
	\item \texttt{metadata\_target}: btree writes (filesystem-level only)
	\item \texttt{background\_target}: If set, user data (not metadata) will
		be moved to this target in the background
	\item\texttt{promote\_target}: If set, a cached copy will be added to
		this target on read, if none exists
\end{description}

\subsubsection{Caching and durability}

Devices can be configured for tiered storage using the target options and
per-device durability settings. The key concept is \emph{durability}: each
device has a durability value (default 1) that controls how many replicas a
copy on that device counts for.

Setting \texttt{durability=0} on a device makes it a pure cache---copies
there do not count toward replication requirements and are evicted in LRU
order when space is needed. Setting \texttt{durability=2} on a hardware RAID
device tells bcachefs the device already has internal redundancy, so one copy
there counts as two replicas.

\paragraph{Tiered storage with replication}

Replication interacts with targets naturally: the filesystem satisfies the
configured \texttt{data\_replicas} count using durability-weighted copies.
Foreground writes place all of the required durable copies on the foreground
target; the reconcile subsystem then moves them to the background target. For
example, with \texttt{data\_replicas=2}, \texttt{foreground\_target=ssd}, and
\texttt{background\_target=hdd}:

\begin{itemize}
\item Foreground writes place both durable copies on the SSD target, assuming
  the SSD devices have enough total durability for \texttt{data\_replicas}
\item Reconcile then moves those copies to the HDD target---both replicas end
  up on the HDD target, not split one per tier
\item The SSD retains a copy only as a cache: set \texttt{promote\_target=ssd}
  to keep a cached copy there for fast reads (added on read, not counted as a
  replica, and evicted in LRU order under space pressure), or set
  \texttt{durability=0} on the SSD devices to use them as a pure write cache
\end{itemize}

See \hyperref[sec:devices]{Device management} in subsystem details for
writeback vs.\ writearound caching configurations and device lifecycle.

\subsubsection{Splitbrain detection}

When assembling a multi-device filesystem, bcachefs detects devices that have
diverged---been mounted and modified independently. Every superblock records a
sequence number (\texttt{seq}) and write timestamp (\texttt{write\_time});
devices with the same sequence but different write times have diverged.

Since version 1.4, superblocks also track the last known sequence of every
other member device---a \emph{vector clock}. A vector clock is a
distributed-systems technique where each participant maintains a vector of
counters, one per participant. Each device's superblock records not just its
own sequence number, but the last sequence number it observed for every other
device. When devices are reassembled, bcachefs compares each device's actual
sequence against what the other devices expected it to be: if device~A's
actual sequence exceeds what device~B recorded for it, then A was modified
independently while B was offline. This catches divergence even in cases where
simple timestamp comparison would not (e.g.\ clock skew, or writes that happen
to produce the same sequence but different content).

The two common causes of splitbrain are:

\begin{itemize}
	\item \textbf{Partial mount of replicated filesystem}: Mounting a subset
		of devices from a replicated filesystem (\texttt{replicas >= 2})
		with \texttt{-o degraded=very} allows writes to proceed despite
		missing devices. If the ``missing'' device is later reconnected,
		it now has stale metadata while the mounted devices have advanced.

	\item \textbf{External snapshots}: LVM, SAN, or VM snapshots create
		block devices with valid bcachefs superblocks sharing the same
		filesystem UUID. If both original and snapshot are visible
		during assembly, or if the snapshot is mounted read-write and
		later reassembled with the original, the devices have diverged.
\end{itemize}

When divergence is detected, bcachefs excludes the stale device. Once devices
have genuinely diverged with writes on both sides, recovery is impractical -
resolving differences key-by-key across all btrees is not feasible. The
\texttt{-o no\_splitbrain\_check} option includes divergent devices anyway,
but is only safe when no actual modifications occurred on the excluded device.

\subsection{Reflink}

bcachefs supports reflink (\texttt{cp --reflink}, \texttt{FICLONE} ioctl),
creating copies that share underlying storage. The original extent is moved to
a separate reflink btree with a reference count; reads through a reflinked
extent require two btree lookups instead of one. Writing new data over a
reflinked range works through normal COW --- the refcount is decremented, and
when it reaches zero the shared extent and its disk space are freed.

Reflink is currently a one-way transformation: once an extent becomes indirect,
it stays indirect even when only one reference remains. This means a file that
was once reflinked permanently pays the double-lookup cost on reads.

IO option propagation (compression, checksums, replicas) for shared extents is
controlled by a security boundary: only the source file's options can propagate
to the shared data, preventing a reflink copy from degrading another user's data
protection settings.

See \hyperref[sec:reflink]{Reflink} in subsystem details for the full
on-disk representation, lifecycle mechanics, and snapshot interaction.

\subsection{Subvolumes and snapshots}

bcachefs supports subvolumes and snapshots with a similar userspace interface as
btrfs. A new subvolume may be created empty, or it may be created as a snapshot
of another subvolume. Snapshots are writeable by default and may be snapshotted
again, creating a tree of snapshots; they can also be created as read-only with
the \texttt{--read-only} (or \texttt{-r}) flag.

Snapshots are very cheap to create: they're not based on cloning of COW btrees
as with btrfs, but instead are based on versioning of individual keys in the
btrees. Many thousands or millions of snapshots can be created, with the only
limitation being disk space.

The following subcommands exist for managing subvolumes and snapshots:
\begin{itemize}
	\item \texttt{bcachefs subvolume create}: Create a new, empty subvolume
	\item \texttt{bcachefs subvolume delete}: Delete an existing subvolume
		or snapshot
	\item \texttt{bcachefs subvolume snapshot}: Create a snapshot of an
		existing subvolume
\end{itemize}

A subvolume can also be deleted with a normal rmdir after deleting all the
contents, as with \texttt{rm -rf}. Still to be implemented: recursive snapshot
creation and a method for recursively listing subvolumes.

See \hyperref[sec:snapshots]{Subvolumes and snapshots} in subsystem details
for the full architecture, key visibility model, and snapshot lifecycle.

\subsection{Quotas}

bcachefs supports conventional user/group/project quotas. Quotas do not
currently apply to snapshot subvolumes, because if a file changes ownership in
the snapshot it would be ambiguous as to what quota data within that file
should be charged to. Quota accounting resolves every inode through the master
subvolume only, charging sectors and inode counts to the uid/gid/project
recorded there. Writes to snapshot subvolumes bypass quota enforcement entirely.
If there is no master subvolume for a snapshot tree, quota accounting for that
tree is skipped.

When a directory has a project ID set it is inherited automatically by
descendants on creation and rename. When renaming a directory would cause the
project ID to change we return -EXDEV so that the move is done file by file, so
that the project ID is propagated correctly to descendants - thus, project
quotas can be used as subdirectory quotas.

\subsection{32-bit inodes}

The \texttt{inodes\_32bit} option restricts new inode numbers to 32 bits. This
is needed for legacy NFS compatibility (NFSv2 and some older NFSv3
implementations use 32-bit inode numbers), for 32-bit userspace that may
truncate 64-bit inode numbers, and for Wine/Proton/Steam which run 32-bit
Windows applications. This option can be set per-directory, affecting only
files created within that directory tree.

Note that \texttt{inodes\_32bit} silently disables
\texttt{shard\_inode\_numbers\_bits}: there are too few bits in the 32-bit
space to make CPU-based sharding practical.

\subsection{Casefolding}

bcachefs supports case-insensitive filenames via the Linux kernel's Unicode
subsystem (\texttt{CONFIG\_UNICODE}). The \texttt{casefold} option is a
per-directory property that is inherited by newly created subdirectories.

\textbf{Important}: Casefolding can only be enabled on an empty directory.
Existing entries cannot be converted because they would need to be rehashed
with case-folded names. On disk, casefolded directories store both the original
filename and its case-folded form; the case-folded form is used for lookups
while the original is preserved for display.

The Unicode version used for case folding is hardcoded to Unicode 12.1.0. The
kernel will fail to mount (or use an older version correctly) if built with
different Unicode tables.

\subsection{Migration}

\texttt{bcachefs migrate} converts an existing filesystem (ext4, XFS, and
others) to bcachefs in place, without reformatting or copying data to a
separate device. The original data is preserved and incorporated into the new
bcachefs filesystem. After migration, \texttt{bcachefs migrate-superblock}
writes a standard superblock so the filesystem can be mounted normally.
Migration from btrfs is not currently supported because its FIEMAP
implementation does not report which device an extent resides on.

\subsection{Images}

\texttt{bcachefs image create} builds a bcachefs filesystem image from a
directory tree, suitable for embedding in OS images, containers, or VM
disks. \texttt{bcachefs image update} incrementally updates an existing image
to reflect changes in the source directory, minimizing the amount of data
rewritten.

\section{Compatibility and versioning}

The on-disk format gains features from release to release, with two
compatibility guarantees that keep upgrading --- and going back --- predictable.

Forwards: a newer release can always mount a filesystem an older release made.
It upgrades the format while mounting, by running the recovery passes that bring
the metadata up to date --- automatically, with nothing asked of you. It is
worth knowing this happens: the format is updated in place, so once you mount
with a newer release the filesystem is on the newer format. (Support for the
very oldest formats is eventually dropped, so a very large version gap may need
an intermediate release.)

Backwards, by default: an older release can still mount and use a filesystem a
newer release has touched. The exception is an \emph{incompatible} upgrade, which
you ask for explicitly with \texttt{version\_upgrade=incompatible}; it turns on
newer format features older releases cannot read, and so is a one-way step ---
until you run \texttt{bcachefs downgrade}, which runs the recovery passes that
convert the filesystem back to a form an older release accepts.

So: upgrading is automatic and, by default, never costs you the ability to go
back; only an explicit incompatible upgrade does, and \texttt{bcachefs downgrade}
undoes even that.

\section{Management}

\subsection{Formatting}

To format a new bcachefs filesystem use the subcommand \texttt{bcachefs
format}, or \texttt{mkfs.bcachefs}. All persistent filesystem-wide options can
be specified at format time. For an example of a multi device filesystem with
compression, encryption, replication and writeback caching:
\begin{quote} \begin{verbatim}
bcachefs format --compression=lz4               \
                --encrypted                     \
                --replicas=2                    \
                --label=ssd.ssd1 /dev/sda       \
                --label=ssd.ssd2 /dev/sdb       \
                --label=hdd.hdd1 /dev/sdc       \
                --label=hdd.hdd2 /dev/sdd       \
                --label=hdd.hdd3 /dev/sde       \
                --label=hdd.hdd4 /dev/sdf       \
                --foreground_target=ssd	        \
                --promote_target=ssd            \
                --background_target=hdd
\end{verbatim} \end{quote}

In this configuration, data is written to the SSD group first and migrated to
the HDD group in the background. With \texttt{-{}-replicas=2}, the SSD and HDD
copies each count as one replica (since all devices default to
\texttt{durability=1}), satisfying the replication requirement across tiers. To
make the SSDs a pure write cache instead, add \texttt{-{}-durability=0} to the
SSD devices so their copies do not count toward replication.

\subsection{Mounting}

To mount a multi device filesystem, there are two options. You can specify all
component devices, separated by colons, e.g.
\begin{quote} \begin{verbatim}
mount -t bcachefs /dev/sda:/dev/sdb:/dev/sdc /mnt
\end{verbatim} \end{quote}
Or, use the mount.bcachefs tool to mount by filesystem UUID or label:
\begin{quote} \begin{verbatim}
mount -t bcachefs UUID=<uuid> /mnt
mount -t bcachefs LABEL=<label> /mnt
\end{verbatim} \end{quote}

No special handling is needed for recovering from unclean shutdown. Journal
replay happens automatically, and diagnostic messages in the dmesg log will
indicate whether recovery was from clean or unclean shutdown.

The \texttt{-o degraded} option will allow a filesystem to be mounted without
all the devices, but will fail if data would be missing. The
\texttt{-o degraded=very} option can be used to attempt mounting when data
would be missing.

The \texttt{-o verbose} option enables additional log output during the mount
process.

\subsubsection{Mounting at boot}

At boot the system may attempt to mount a multi device filesystem before udev
has finished probing all of its component devices. Since mounting normally
requires every device to be present, this is a race: the mount can fail
reporting missing devices even though they are all attached, simply because
some had not been probed yet when the mount was attempted.

bcachefs-tools ships a templated systemd service to close this race,
\texttt{bcachefs-wait-devices@.service}, which runs \texttt{bcachefs
wait-devices UUID=<uuid>}. It uses udev to watch the block subsystem and does
not complete until every device belonging to the filesystem has been
initialized. Make the mount depend on it by adding to the fstab entry's mount
options
\begin{quote} \begin{verbatim}
x-systemd.requires=bcachefs-wait-devices@<uuid>.service
\end{verbatim} \end{quote}
where \texttt{<uuid>} is the filesystem UUID---the \texttt{@} instance is passed
through to \texttt{wait-devices} as the UUID to wait for. systemd then holds the
mount until all of the filesystem's devices have appeared.

Because the service waits for \emph{every} member device, it should not be
combined with \texttt{-o degraded} when a device is permanently missing: the
service would never complete and the mount would time out. Use the wait service
for the normal, all-devices-present case, and \texttt{degraded} without the
requirement when intentionally mounting a filesystem that is missing a device.

\subsection{Monitoring}

Three commands provide a comprehensive view of filesystem health and
performance:

\begin{description}
\item[\texttt{bcachefs fs usage}] Displays filesystem space usage broken out
  by data type (superblock, journal, btree, data, cached data, parity), by
  which sets of devices extents are replicated across, and per-device usage
  including fragmentation from partially-used buckets. This is the primary
  tool for understanding where space is going and how data is distributed
  across devices.

\item[\texttt{bcachefs fs top}] Real-time display of btree operations by
  process, similar to \texttt{top(1)} for CPU usage. Shows which processes are
  generating btree reads, writes, and transaction restarts. This is the
  primary tool for diagnosing metadata performance problems---if the
  filesystem feels slow, \texttt{fs top} will show what's generating the
  load. Every counter displayed by \texttt{fs top} has a corresponding
  tracepoint, so once you identify which operation is hot, you can drill
  down into individual events with \texttt{perf trace} or \texttt{bpftrace}.

\item[\texttt{bcachefs fs timestats}] Real-time display of operation latency
  statistics: min, max, mean, standard deviation, and EWMA for both duration
  and inter-arrival time. Covers all major IO and internal operations (data
  reads, data writes, btree reads, btree splits, journal flushes, copygc,
  etc.). This is the primary tool for understanding latency---if an
  operation is slow, timestats will show which stage of the pipeline is
  responsible.
\end{description}

For deeper investigation, sysfs
(\texttt{/sys/fs/bcachefs/<uuid>/}) provides per-option tunables, persistent
counters, and internal state; debugfs
(\texttt{/sys/kernel/debug/bcachefs/<uuid>/}) provides btree contents,
transaction state, lock contention, and journal pin diagnostics. See
\S\ref{sec:debugging} for details.

\subsection{Journal tuning}

The journal has tunables that affect filesystem performance; see
\hyperref[sec:journal]{Journal} in subsystem details for the full architecture.
The key tunables:

\begin{description}
\item[\texttt{journal\_flush\_delay}] Milliseconds before auto-flushing
  (default 1000). Sync and fsync force immediate flushes.
\item[\texttt{journal\_flush\_disabled}] Disables sync/fsync flushes. The
  \texttt{journal\_flush\_delay} timer still applies, so at most one second of
  work is lost on crash. Useful on workstations, less appropriate on servers.
\item[\texttt{journal\_reclaim\_delay}] Background reclaim interval (default
  100\,ms).
\end{description}

The journal should be sized so that bursts of activity do not fill it; a
larger journal also allows larger batched btree writes. Use
\texttt{bcachefs device resize-journal} to grow the journal online. Current
utilization is visible via
\texttt{/sys/fs/bcachefs/<uuid>/internal/journal\_debug}.

\subsection{Device management}

Devices can be added, removed, resized, and taken online/offline at any time
without unmounting. The \texttt{bcachefs device} subcommands handle all device
lifecycle operations; see \hyperref[sec:devices]{Device management} in
subsystem details for the full treatment of device states, durability,
caching, degraded mode, and error tracking.

\subsection{Data management}

\subsubsection{Reconcile}

The reconcile subsystem (formerly ``rebalance'') runs in the background to
ensure all data and metadata matches configured IO path options: replication
count (\texttt{data\_replicas}, \texttt{metadata\_replicas}), compression,
checksum type, erasure coding, and device targets (\texttt{foreground\_target},
\texttt{metadata\_target}, etc.). When options change or devices are added or
removed, reconcile propagates these changes to existing data and metadata by
moving, rewriting, or recompressing extents and btree nodes.

Reconcile uses 6 dedicated btrees for work tracking:
\texttt{reconcile\_work} and \texttt{reconcile\_hipri} for normal and
high-priority logical work, \texttt{reconcile\_pending} for work that failed
due to insufficient space or devices, \texttt{reconcile\_scan} for option
change propagation via scan cookies, and physical variants
(\texttt{reconcile\_work\_phys}, \texttt{reconcile\_hipri\_phys}) that track
work by device LBA for efficient processing on rotational devices.

Work enters the system when filesystem or inode options change (triggering a
scan), or when extent triggers detect a mismatch between current data placement
and desired options. The reconcile thread processes work in priority order:
high-priority metadata first (under-replicated or evacuating), then high-priority
data, then normal metadata (moving stray metadata to \texttt{metadata\_target}),
then normal data, then pending retries (only after device configuration changes).
The \texttt{reconcile\_pending} btree solves the long-standing ``rebalance
spinning'' problem, where the old rebalance thread would burn CPU retrying
moves that could not complete (e.g.\ because the target was full). A practical
consequence: you can create a single-device filesystem with
\texttt{replicas=2}, and all data will be marked as degraded with 1x
availability; when a second device is added, reconcile automatically
replicates everything to 2x without further user intervention.

The \texttt{bcachefs reconcile status} command shows current reconcile
progress, and \texttt{bcachefs reconcile wait} blocks until reconcile
completes for specified work types. The \texttt{reconcile\_enabled} and
\texttt{reconcile\_on\_ac\_only} options control when reconcile runs.

\subsubsection{Scrub}

Scrub reads all data on a running filesystem and verifies checksums, detecting
silent data corruption (bitrot). When a checksum mismatch is found and a valid
redundant copy exists (from replication or erasure coding), the corrupted copy
is automatically repaired by rewriting it from the good copy. Scrub walks data
in physical (LBA) order using \hyperref[sec:backpointers]{backpointers}, which is efficient for rotational
devices and avoids the random access pattern that would result from walking the
logical extent tree.

Scrub can be run on a specific device or on all devices in the filesystem.
Progress is reported via sysfs and can be monitored with
\texttt{bcachefs data scrub}. Affected file paths are currently logged to the
kernel log only.

\section{Repair and recovery}

bcachefs is designed to detect, report, and repair problems at every level:
inline during normal IO, in the background via scheduled recovery passes, and
through explicit administrator-driven tools. This section covers all of these
mechanisms.

\subsection{Online self-healing}

Bcachefs classifies metadata errors into three tiers based on severity and
whether a safe automatic repair is known:

\begin{description}
\item[Autofix errors (\texttt{FSCK\_AUTOFIX})] Inconsistencies with a
  well-understood, always-safe repair that can be applied immediately inline
  during normal operation, without interrupting I/O or requiring a recovery
  pass. An example is a missing journal entry replica marker: the repair is
  simply to write the correct marker, with no data at risk.

\item[Fixable errors (\texttt{FSCK\_CAN\_FIX})] Inconsistencies that have a
  known repair procedure but require a dedicated recovery pass to execute
  safely. Under \texttt{errors=fix\_safe} (the default since version 1.19),
  detecting such an error schedules the relevant recovery pass to run online
  in the background. The pass runs concurrently with normal filesystem
  operation without taking the filesystem offline, provided it carries the
  \texttt{PASS\_ONLINE} flag. Errors detected during recovery that would
  rewind the recovery sequence are handled specially: recovery restarts from
  the affected pass rather than aborting.

\item[Inconsistent errors] Structural corruption for which no safe automatic
  repair is known, or where continuing operation risks further data loss. These
  trigger emergency read-only regardless of the \texttt{errors} setting (except
  \texttt{errors=continue}, which logs the error and proceeds). The filesystem
  logs a message directing the administrator to run fsck. If a repair procedure
  is subsequently developed and deemed safe, the error can be reclassified to
  \texttt{FSCK\_CAN\_FIX} in a future kernel release.
\end{description}

The \texttt{errors} mount option controls the policy for the second and third
tiers; see \S\ref{sec:error-actions}. The default \texttt{fix\_safe} enables
background self-repair for fixable errors while taking the conservative
read-only path for inconsistent ones.

\textbf{Read path self-healing}: When a read detects a checksum mismatch and
a valid replica exists, the corrupted copy is automatically repaired by
rewriting it from the good copy. This happens transparently on every read
without administrator intervention. If erasure coding is available, data can
be reconstructed from parity even with no good replica.

\textbf{Error accounting}: Every error detection event is recorded as a
persistent count in the superblock's errors section. These counts survive
across reboots and accumulate over the filesystem's lifetime, providing an
audit trail of what has been detected and repaired. The counts are visible in
\texttt{bcachefs show-super} output and via the sysfs \texttt{errors} file.

\subsection{Fsck}

The \texttt{bcachefs fsck} subcommand (also available as \texttt{fsck.bcachefs})
can run fsck in userspace (\texttt{-K}), in the kernel at mount time
(\texttt{-k}), or online on a mounted filesystem, selecting the mode
automatically or via command line options. The in-kernel offline mode is useful
when an older bcachefs-tools binary cannot handle the filesystem version. In
all cases the exact same fsck implementation runs---the recovery passes
described in the introduction. Running in userspace allows stopping with
ctrl-c and prompting for fixes; \texttt{bcachefs fsck -y} auto-fixes,
\texttt{-n} runs read-only.

Online fsck runs the validation passes marked as safe for concurrent use
(\texttt{PASS\_ONLINE}) on a mounted, read-write filesystem. Combined with
\texttt{errors=fix\_safe}, this means most metadata inconsistencies are
detected and repaired automatically without administrator intervention.

\subsubsection{Safe inspection with nochanges}

The \texttt{-o nochanges} mount option disallows any writes to the underlying
devices, pinning dirty data in memory if journal replay is necessary. This is
the primary tool for safe inspection of a damaged filesystem: it allows
recovery to run and metadata to be read without risking further damage.

\texttt{bcachefs fsck -n} implies \texttt{-o nochanges}. To test a repair
before committing:
\texttt{-o ro,nochanges,fsck,fix\_errors} runs the full repair in memory
without writing to disk, allowing the filesystem to be mounted and its
contents verified before repeating without \texttt{nochanges}.

\subsection{Journal rewind}

Journal rewind rolls back filesystem metadata to an earlier point in time by
undoing committed transactions. Each transaction records both the new key and
the old key it replaced; rewind replays the old keys to restore the previous
state, then triggers a full fsck to repair allocation metadata. See
\S\ref{sec:journal} for the underlying mechanism.

Rewind serves two purposes: automatic recovery from bad device flushes, and
manual disaster recovery.

\subsubsection{Automatic rewind via journal scrub}

The \hyperref[opt:scrub-recent-journal-entries]{\texttt{scrub\_recent\_journal\_entries}}
mount option can trigger journal rewind automatically. During mount, it verifies
the data referenced by recent journal entries between flush boundaries. Each
replica of each extent is read and checksum-verified independently. The scrub
walks flush ranges newest-first and stops after two consecutive ranges with no
errors, or when it reaches the time limit set by
\texttt{scrub\_journal\_max\_rewind\_secs} (default 10 seconds).

The response depends on whether the filesystem has redundancy:

\begin{description}
\item[No redundancy (single device or no replication):] If any data fails
  checksum verification, the journal is rewound to the last flush boundary
  before the corrupted entry, discarding the metadata that pointed to bad data.
  Recovery then continues with a full fsck from the rewound state.
\item[With replication:] When only some replicas of an extent are bad but at
  least one is good, there is no need to rewind---the good replica can serve
  reads. Instead, the bad extents are stashed in a repair list and fixed
  immediately after the filesystem goes read-write, by rewriting the corrupted
  replicas from the good copies. This avoids discarding valid metadata when the
  data itself is still recoverable. Rewind only happens if \emph{all} replicas
  of an extent are corrupted.
\end{description}

This is the primary defense against drives with buggy cache flush handling: the
filesystem detects at mount time that acknowledged writes never reached stable
storage, and recovers automatically without administrator intervention.
Offending devices are identified and a per-device \texttt{flush\_errors}
counter in the superblock is incremented, providing a persistent record visible
in \texttt{bcachefs show-super} output.

\subsubsection{Manual rewind}

The \hyperref[opt:journal-rewind]{\texttt{journal\_rewind}} mount option
accepts a journal sequence number and rewinds to that point. This is a disaster
recovery tool for reverting damage from a bug that irreversibly corrupted or
deleted metadata, or (combined with
\hyperref[opt:nochanges]{\texttt{-o nochanges}}) for attempting to recover
recently deleted files.

Use \texttt{bcachefs list\_journal} to identify the sequence number to rewind
to. The \texttt{journal\_transaction\_names} option (on by default) must have
been enabled when the entries were written---it controls whether the old keys
needed for rewind are stored in the journal.

\subsubsection{Limitations}

Since rewind operates on metadata only, data in buckets that have been reused
(generation incremented and overwritten with new data) is unrecoverable, and
\texttt{nocow} data cannot be rolled back at all. The journal records allocator
updates, so \texttt{list\_journal -b alloc} can identify when specific buckets
were recycled. The rewind window is bounded by the journal size: only sequences
between \texttt{last\_seq} (oldest entry still on disk) and the current
sequence are available, so a larger journal provides a deeper rewind history.

\subsection{Disaster recovery}

bcachefs is designed to recover from catastrophic metadata damage---not just
the routine inconsistencies that fsck handles, but situations where entire
btrees are lost or corrupted beyond normal repair. The key insight is that
the filesystem's data is recoverable as long as the leaf nodes of the extents
and dirents btrees survive, because those two btrees contain the essential
mappings: which data belongs to which file, and which files exist in which
directories. Everything else can be reconstructed from them.

In practical terms: if every btree except extents and dirents is destroyed,
recovery can still place every file's data in its correct directory with its
correct name. File sizes may be slightly off if inode metadata is lost
(inodes record the precise file size, while extents only record block-aligned
ranges), and permissions, ownership, and timestamps will be lost with the
inodes, but the actual data---the bytes the user cares about---will be intact
and in the right place. This is a much stronger recovery guarantee than most
filesystems provide.

\subsubsection{Reconstruction passes}

When normal recovery detects structural damage, it automatically schedules
heavier reconstruction passes:

\begin{description}
\item[\texttt{scan\_for\_btree\_nodes} + \texttt{check\_topology}] These two
  passes work together. When \texttt{check\_topology} finds missing or
  unreachable btree nodes, it schedules \texttt{scan\_for\_btree\_nodes}, which
  physically scans every device for valid btree node headers (identified by
  magic number), deduplicates replicas, and builds a table of recovered nodes.
  Then \texttt{check\_topology} runs again, using the scan results to
  reconstruct missing interior nodes and reconnect orphaned subtrees. Together
  they can recover a btree even when the root node and multiple levels of
  interior nodes are lost---as long as the leaf nodes are intact. This reads
  every sector of every device and may take considerable time on large
  filesystems.

\item[\texttt{reconstruct\_snapshots}] Scheduled when the snapshots btree
  itself is lost, or when references to missing snapshots cannot be resolved
  by deletion alone. Scans all snapshot-aware btrees to find which snapshot
  IDs are actually in use, then reconstructs the snapshot tree structure from
  that evidence.

\item[\texttt{check\_allocations}] Full allocation rebuild: walks all btrees
  marking which buckets are referenced, then compares against the alloc btree
  to repair data types, sector counts, and stripe references. This is the
  pass that makes it possible to lose the entire alloc btree and recover---the
  allocation state is fully derivable from the extent and backpointer btrees.
\end{description}

These passes are not part of normal recovery and are only triggered when
damage is detected. They run automatically---no special flags or manual
intervention is required. After reconstruction, the filesystem schedules
a full fsck to validate the rebuilt metadata.

\subsubsection{Superblock recovery}

The superblock itself has multiple layers of redundancy. Up to 61 backup
copies can be stored on each device, and every device in the filesystem has
a complete copy of the superblock. The \texttt{bcachefs recover-super} command
can reconstruct a completely overwritten primary superblock from backup copies
on the same device or from any other device in the filesystem.

\subsubsection{Drive firmware failures}

Storage devices almost always have volatile write caches: a write is sent and
acknowledged, but not actually written to stable media until a separate flush
operation. Unfortunately, not all devices handle flushes reliably; this is
particularly known to be a problem on cheaper devices (including at the
controller level), and has historically been a major (potentially disastrous)
problem for COW filesystems.

bcachefs mitigates this at multiple levels:

\begin{itemize}
\item \textbf{Checksumming} detects the corruption when the data is eventually
  read, preventing silent data corruption from propagating to applications.
\item \textbf{Replication} provides a correct copy to serve reads from and
  repair the corrupted copy automatically.
\item \textbf{Journal scrub} (\hyperref[opt:scrub-recent-journal-entries]{\texttt{scrub\_recent\_journal\_entries}}) detects
  the problem at mount time by verifying recently written data from the
  previous mount. If problems are found, the bad flush will be logged and
  recorded in the superblock entry for the faulty member device, and then the
  data will be immediately repaired from other replicas (on a replicated
  filesystem), or if there are no good replicas the entire filesystem will be
  rolled back to the previous consistent state (with a default maximum rollback
  of 10 seconds).
\end{itemize}

With journal scrub enabled, even a single-device filesystem without replication
recovers cleanly from a flush-lying drive: the corrupted entries are detected
at mount time and the journal is rewound, losing only the unflushed data (which
would have been lost in a legitimate power failure anyway) rather than leaving
files permanently unreadable. With replication, recovery is even better: the
metadata is preserved and only the corrupted replicas are rewritten.

\subsubsection{Diagnostic tools}

When investigating damage before attempting repair:

\begin{description}
\item[\texttt{bcachefs fsck -n}] Runs all recovery and repair passes without
  writing to disk (\texttt{-o nochanges}). Shows what would be repaired.
\item[\texttt{bcachefs list}] Dumps btree contents for offline inspection.
  Useful for examining what data survived in specific btrees.
\item[\texttt{bcachefs list\_journal}] Lists journal contents with sequence
  numbers, transaction names, and per-btree filtering. Essential for
  understanding what happened before a crash and for identifying journal
  rewind targets.
\item[\texttt{bcachefs dump}] Exports all metadata as qcow2 images for
  offline analysis, without including user data. For encrypted filesystems,
  the passphrase must first be removed with \texttt{bcachefs remove-passphrase}.
\end{description}

\subsubsection{Extent poisoning and data recovery}

When a read detects a checksum mismatch and no valid replica exists to serve
the read from, the extent is marked \emph{poisoned}. Subsequent reads of a
poisoned extent return an error immediately without going to disk, preventing
applications from silently consuming corrupt data.

Poisoning is persistent: it survives remounts and data moves. When a poisoned
extent is moved (by reconcile, copygc, or tiering), the filesystem must
compute a new checksum over the data at its new location. Because the data is
corrupt, the new checksum is valid for the corrupt data. The poison flag is how
the filesystem remembers that the data was known-bad before the move---without
it, the new checksum would pass and the corruption would become invisible.

\paragraph{Inspecting poisoned data.}
The \texttt{bcachefs data-read} command performs an O\_DIRECT read with
extended error reporting. It reports errors via a bitmask (checksum, IO,
decompression, erasure coding reconstruction) and detailed kernel error
messages.

\begin{verbatim}
# Read with error diagnostics
bcachefs data-read /path/to/file --offset 0 --len 4096

# Read poisoned data (bypasses poison check)
bcachefs data-read /path/to/file --no-poison-check

# Save recovered data to a file
bcachefs data-read /path/to/file --no-poison-check --output recovered.bin
\end{verbatim}

With \texttt{--no-poison-check}, the command bypasses the poison check and
returns whatever data is on disk. The errors bitmask will include
\texttt{checksum} to indicate that the extent was poisoned.

\paragraph{Clearing poison flags.}
If the data has been verified as acceptable (or restored from backup), the
\texttt{bcachefs unpoison} command clears the poison flag:

\begin{verbatim}
bcachefs unpoison /path/to/file --yes-i-understand
\end{verbatim}

The \texttt{--yes-i-understand} flag is required because unpoisoning removes
the only indication that the data is bad. After unpoisoning, normal reads
return the data without errors, even if it is corrupt. This is appropriate
when the data is usable despite corruption (e.g.\ a media file with minor
artifacts) or when the file has been restored from backup.

\section{Command reference}

The \texttt{bcachefs} tool provides filesystem management through grouped
subcommands. Run \texttt{bcachefs} with no arguments for a command summary,
or \texttt{bcachefs <command> --help} for detailed usage.

\bchdoc{cli-reference}

\section{Options}
\label{sec:options}

Most bcachefs options can be set filesystem wide, and a significant subset can
also be set on inodes (files and directories), overriding the global defaults.
Filesystem wide options may be set when formatting, when mounting, or at runtime
via \texttt{/sys/fs/bcachefs/<uuid>/options/}. When set at runtime via sysfs,
the persistent options in the superblock are updated as well; when options are
passed as mount parameters the persistent options are unmodified. Additionally
some of the filesystem wide options can be set via \texttt{bcachefs set-fs-option}

\subsection{Options reference}

\bchdoc{opts-table}

\subsection{File and directory options}

A subset of filesystem options can be set on individual files and directories,
overriding the filesystem-wide defaults. The per-inode options are:
\texttt{data\_checksum}, \texttt{compression}, \texttt{background\_compression},
\texttt{data\_replicas}, \texttt{foreground\_target},
\texttt{background\_target}, \texttt{promote\_target},
\texttt{erasure\_code}, \texttt{nocow}, \texttt{casefold}, and
\texttt{inodes\_32bit}. Note that \texttt{casefold} and \texttt{inodes\_32bit}
can only be toggled on empty directories.

Each inode stores its own copy of these options so that reads never need to
walk parent directories. Each option has a separate ``defined'' flag that
records whether the option was explicitly set on this inode; options without
this flag are propagated from the parent directory on inode creation and rename.
When renaming a directory would cause inherited attributes to change we fail the
rename with -EXDEV, causing userspace to do the rename file by file so that
inherited attributes stay consistent.

Inode options are available as extended attributes. The options that have been
explicitly set are available under the \texttt{bcachefs} namespace, and the
effective options (explicitly set and inherited options) are available under the
\texttt{bcachefs\_effective} namespace. Examples of listing options with the
getfattr command:

\begin{quote} \begin{verbatim}
$ getfattr -d -m '^bcachefs\.' filename
$ getfattr -d -m '^bcachefs_effective\.' filename
\end{verbatim} \end{quote}

Options may be set via the extended attribute interface, but it is preferable to
use the \texttt{bcachefs set-file-option} command as it will correctly
propagate options recursively.

\subsection{Inode number sharding}

On systems with many CPUs, inode number allocation can become a contention
point: every file creation needs a new inode number, and a single global cursor
serializes all CPUs. The \texttt{shard\_inode\_numbers\_bits} option (0-8,
set at format time to the log2 of the number of online CPUs if not specified)
partitions the inode number space by CPU ID, giving each CPU its own
allocation range. This eliminates cross-CPU contention at the cost of
non-sequential inode numbers.

\subsection{Error actions}
\label{sec:error-actions}
The \texttt{errors} option controls how the filesystem responds to metadata
inconsistencies. Valid error actions are:
\bchdoc{error-actions}

\subsection{Checksum types}
Valid checksum types are:
\bchdoc{csum-opts}

\subsection{Compression types}
Valid compression types are:
\bchdoc{compression-opts}

\subsection{String hash types}
Valid hash types for string hash tables are:
\bchdoc{str-hash-opts}

\section{Debugging tools}
\label{sec:debugging}

\subsection{Sysfs interface}

Mounted filesystems are available in sysfs at \texttt{/sys/fs/bcachefs/<uuid>/}
with various options, performance counters and internal debugging aids.

\subsubsection{Options}

Filesystem options may be viewed and changed via \\
\texttt{/sys/fs/bcachefs/<uuid>/options/}, and settings changed via sysfs will
be persistently changed in the superblock as well.

\subsubsection{Time stats}
\label{sec:timestats}

bcachefs tracks the latency and frequency of various operations and events, with
quantiles for latency/duration in the
\texttt{/sys/fs/bcachefs/<uuid>/time\_stats/} directory.

\bchdoc{time-stats}

\subsubsection{Persistent counters}
\label{sec:counters}

bcachefs tracks persistent counters in the superblock that survive across mounts.
These are available in \texttt{/sys/fs/bcachefs/<uuid>/counters/} and count
either events or sectors depending on type.

\bchdoc{counters}

\subsubsection{Internals}

\begin{description}
	\item \texttt{btree\_cache} \\
		Shows information on the btree node cache: number of cached
		nodes, number of dirty nodes, and whether the cannibalize lock
		(for reclaiming cached nodes to allocate new nodes) is held.

	\item \texttt{btree\_key\_cache} \\
		Prints information on the btree key cache: number of freed keys
		(which must wait for an SRCU barrier to complete before being
		freed), number of cached keys, and number of dirty keys.

	\item \texttt{journal\_debug} \\
		Prints a variety of internal journal state.

	\item \texttt{new\_stripes} \\
		Lists new erasure-coded stripes being created.

	\item \texttt{open\_buckets} \\
		Lists buckets currently being written to, along with data type
		and refcount.

	\item \texttt{io\_timers\_read} \\
	\item \texttt{io\_timers\_write} \\
		Lists outstanding IO timers - timers that wait on total reads or
		writes to the filesystem.

	\item \texttt{trigger\_journal\_commit} \\
		Echoing to this file triggers a journal commit.

	\item \texttt{trigger\_journal\_flush} \\
		Echoing to this file flushes all journal pins (forcing dirty
		btree nodes and key cache entries to be written) and then
		issues a journal meta write.

	\item \texttt{trigger\_gc} \\
		Echoing to this file causes the GC code to recalculate each
		bucket's oldest\_gen field.

	\item \texttt{trigger\_btree\_cache\_shrink} \\
		Echoing to this file prunes the btree node cache.

\end{description}

\subsubsection{Unit and performance tests}

Echoing into \texttt{/sys/fs/bcachefs/<uuid>/perf\_test} runs various low level
btree tests, some intended as unit tests and others as performance tests. The
syntax is
\begin{quote} \begin{verbatim}
	echo <test_name> <nr_iterations> <nr_threads> > perf_test
\end{verbatim} \end{quote}

When complete, the elapsed time will be printed in the dmesg log. The full list
of tests that can be run can be found near the bottom of
\texttt{fs/bcachefs/tests.c}.

\subsection{Debugfs interface}

Various internal debugging files are available under
\texttt{/sys/kernel/debug/bcachefs/<uuid>/}.

The \texttt{btrees/} subdirectory contains one directory per btree (e.g.
\texttt{btrees/extents/}, \texttt{btrees/inodes/}), each with three files:

\begin{description}
	\item \texttt{keys} \\
		Entire btree contents, one key per line.

	\item \texttt{formats} \\
		Information about each btree node: the size of the packed bkey
		format, how full each btree node is, number of packed and
		unpacked keys, and number of nodes and failed nodes in the
		in-memory search trees.

	\item \texttt{bfloat-failed} \\
		For each sorted set of keys in a btree node, we construct a
		binary search tree in eytzinger layout with compressed keys.
		Sometimes we aren't able to construct a correct compressed
		search key, which results in slower lookups; this file lists the
		keys that resulted in these failed nodes.
\end{description}

The remaining files in the debugfs root are:

\begin{description}
	\item \texttt{btree\_transactions} \\
		Lists each running btree transaction that has locks held,
		listing which nodes they have locked and what type of lock, what
		node (if any) the process is blocked attempting to lock, and
		where the btree transaction was invoked from.

	\item \texttt{btree\_updates} \\
		Lists outstanding interior btree updates: the mode (nothing
		updated yet, or updated a btree node, or wrote a new btree root,
		or was reparented by another btree update), whether its new
		btree nodes have finished writing, its embedded closure's
		refcount (while nonzero, the btree update is still waiting), and
		the pinned journal sequence number.

	\item \texttt{journal\_pins} \\
		Lists what is preventing each journal entry from being
		reclaimed: dirty btree nodes that still reference that
		sequence number, key cache entries waiting to be flushed,
		or in-progress operations that need the entry to remain
		available for crash recovery. Useful for diagnosing journal
		space exhaustion.

	\item \texttt{btree\_deadlock} \\
		Checks for and reports btree lock deadlock cycles among
		transactions.

	\item \texttt{cached\_btree\_nodes} \\
		Lists btree nodes currently in the btree node cache.

	\item \texttt{btree\_transaction\_stats} \\
		Per-transaction-function statistics: lock hold times, restart
		counts, and maximum memory usage.

	\item \texttt{btree\_node\_scan} \\
		Results from the last scan for btree nodes on raw devices,
		used during catastrophic recovery.

	\item \texttt{write\_points} \\
		Lists active write points and the buckets they are currently
		writing to.
\end{description}

\subsection{Listing and dumping filesystem metadata}

\subsubsection{bcachefs show-super}

This subcommand is used for examining and printing bcachefs superblocks. It
takes two optional parameters:
\begin{description}
	\item \texttt{-l}: Print superblock layout, which records the amount of
		space reserved for the superblock and the locations of the
		backup superblocks.
	\item \texttt{-f, --fields=(fields)}: List of superblock sections to
		print, \texttt{all} to print all sections.
\end{description}

\subsubsection{bcachefs list}

This subcommand gives access to the same functionality as the debugfs interface,
listing btree nodes and contents, but for offline filesystems.

\subsubsection{bcachefs list\_journal}

This subcommand lists the contents of the journal for offline filesystems. The
journal is a complete record of every btree update, accounting delta, log
message, and timestamp, ordered by sequence number. Each entry can be traced
back to the transaction function that created it (when
\texttt{journal\_transaction\_names} is enabled). Filtering by btree, sequence
range, or transaction name makes it possible to reconstruct the exact sequence
of events leading to a problem - and to identify the right sequence number for
journal rewind.

\subsubsection{bcachefs dump}

This subcommand can dump all metadata in a filesystem (including multi device
filesystems) as qcow2 images: when encountering issues that \texttt{fsck} can
not recover from and need attention from the developers, this makes it possible
to send the developers only the required metadata. Encrypted filesystems must
first be unlocked with \texttt{bcachefs remove-passphrase}.

\section{Subsystem details}

\subsection{Data paths}

This section describes the end-to-end flow of read and write operations through
the filesystem.

\bchdoc{data-paths}

\subsubsection{Extent checksums and compression}

\bchdoc{extent-checksums}

\subsection{Allocator}

\bchdoc{allocator}

\subsection{Failure domains}
\label{sec:failure-domains}

\bchdoc{failure-domains}

\subsection{Subvolumes and snapshots}
\label{sec:snapshots}

\bchdoc{snapshots}

\subsection{Superblock}
\label{sec:superblock}

The superblock stores all filesystem-wide configuration and per-device
metadata. Each device maintains its own copy, and the copies are kept in sync:
most fields are shared across all devices, with the exception of the device
index and journal bucket locations.

\bchdoc{superblock}

\subsection{Device management}
\label{sec:devices}

\bchdoc{device-management}

\subsection{Journal}
\label{sec:journal}

\bchdoc{journal}

\subsection{Btrees}
\label{sec:btrees}

\bchdoc{btrees}

\subsection{Erasure coding}
\label{sec:erasure-coding}

\bchdoc{erasure-coding}

\section{ioctl interface}

This section documents bcachefs-specific ioctls:

\begin{description}
	\item \texttt{BCH\_IOCTL\_QUERY\_UUID} \\
		Returns the UUID of the filesystem: used to find the sysfs
		directory given a path to a mounted filesystem.

	\item \texttt{BCH\_IOCTL\_FS\_USAGE} \\
		Obsolete. Formerly queried filesystem usage; replaced by
		\texttt{BCH\_IOCTL\_QUERY\_ACCOUNTING}.

	\item \texttt{BCH\_IOCTL\_DEV\_USAGE} \\
		Obsolete. Formerly queried per-device usage; replaced by
		\texttt{BCH\_IOCTL\_QUERY\_ACCOUNTING}.

	\item \texttt{BCH\_IOCTL\_READ\_SUPER} \\
		Returns the filesystem superblock, and optionally the superblock
		for a particular device given that device's index.

	\item \texttt{BCH\_IOCTL\_DISK\_ADD} \\
		Given a path to a device, adds it to a mounted and running
		filesystem. The device must already have a bcachefs superblock;
		options and parameters are read from the new device's superblock
		and added to the member info section of the existing
		filesystem's superblock.

	\item \texttt{BCH\_IOCTL\_DISK\_REMOVE} \\
		Given a path to a device or a device index, attempts to remove
		it from a mounted and running filesystem. This operation
		requires walking the btree to remove all references to this
		device, and may fail if data would become degraded or lost,
		unless appropriate force flags are set.

	\item \texttt{BCH\_IOCTL\_DISK\_ONLINE} \\
		Given a path to a device that is a member of a running
		filesystem (in degraded mode), brings it back online.

	\item \texttt{BCH\_IOCTL\_DISK\_OFFLINE} \\
		Given a path or device index of a device in a multi device
		filesystem, attempts to close it without removing it, so that
		the device may be re-added later and the contents will still be
		available.

	\item \texttt{BCH\_IOCTL\_DISK\_SET\_STATE} \\
		Given a path or device index of a device in a multi device
		filesystem, attempts to set its state to one of read-write,
		read-only, evacuating, or spare. Takes flags to force if the
		filesystem would become degraded.

	\item \texttt{BCH\_IOCTL\_DISK\_GET\_IDX} \\
	\item \texttt{BCH\_IOCTL\_DISK\_RESIZE} \\
	\item \texttt{BCH\_IOCTL\_DISK\_RESIZE\_JOURNAL} \\
	\item \texttt{BCH\_IOCTL\_DATA} \\
		Starts a data job, which walks all data and/or metadata in a
		filesystem, performing some operations on each btree
		node and extent. Returns a file descriptor which can be read
		from to get the current status of the job, and closing the file
		descriptor (i.e. on process exit stops the data job.

	\item \texttt{BCH\_IOCTL\_SUBVOLUME\_CREATE} \\
	\item \texttt{BCH\_IOCTL\_SUBVOLUME\_DESTROY} \\

	\item \texttt{BCH\_IOCTL\_FSCK\_OFFLINE} \\
		Runs offline fsck on an unmounted filesystem via ioctl,
		reporting progress and results through a file descriptor.

	\item \texttt{BCH\_IOCTL\_FSCK\_ONLINE} \\
		Runs online fsck on a mounted filesystem, executing the
		same recovery passes that \texttt{-o fsck} would run at
		mount time.

	\item \texttt{BCH\_IOCTL\_QUERY\_ACCOUNTING} \\
		Queries disk accounting data from the accounting btree:
		per-replica-set usage, per-device usage, compression
		ratios, and other counters. Replaces the obsolete
		\texttt{FS\_USAGE} and \texttt{DEV\_USAGE} ioctls.

	\item \texttt{BCH\_IOCTL\_SUBVOLUME\_LIST} \\
		Lists subvolumes in a filesystem.

	\item \texttt{BCH\_IOCTL\_SNAPSHOT\_TREE} \\
		Queries the full snapshot tree with per-node disk accounting.

	\item \texttt{BCHFS\_IOC\_REINHERIT\_ATTRS} \\
\end{description}

Most disk management ioctls (\texttt{DISK\_ADD}, \texttt{DISK\_REMOVE},
\texttt{DISK\_ONLINE}, \texttt{DISK\_OFFLINE}, \texttt{DISK\_SET\_STATE},
\texttt{DISK\_RESIZE}, \texttt{DISK\_RESIZE\_JOURNAL},
\texttt{SUBVOLUME\_CREATE}, \texttt{SUBVOLUME\_DESTROY}) have v2 variants that
add structured error reporting via an embedded error message buffer.

\section{On disk format}

\subsection{Superblock}

The superblock is the first thing to be read when accessing a bcachefs
filesystem. It is located 4kb from the start of the device, with redundant
copies elsewhere - typically one immediately after the first superblock, and one
at the end of the device.

The \texttt{bch\_sb\_layout} records the amount of space reserved for the
superblock as well as the locations of all the superblocks. It is included with
every superblock, and additionally written 3584 bytes from the start of the
device (512 bytes before the first superblock).

Most of the superblock is identical across each device. The exceptions are the
\texttt{dev\_idx} field, and the journal section which gives the location of the
journal.

The main section of the superblock contains UUIDs, version numbers, number of
devices within the filesystem and device index, block size, filesystem creation
time, and various options and settings. The superblock also has a number of
variable length sections:

\bchdoc{sb-fields}

Several field types have versioned variants (e.g. \texttt{members\_v1}/\texttt{v2},
\texttt{replicas\_v0}/\texttt{replicas}, \texttt{journal}/\texttt{journal\_v2})
for backward compatibility; the kernel reads both old and new formats but writes
only the current version.

\subsection{Journal}

Every journal write (\texttt{struct jset}) contains a list of entries:
\texttt{struct jset\_entry}. Below are listed the various journal entry types.

\bchdoc{jset-entry-types}

\subsection{Btrees}

bcachefs uses 28 separate btrees for different data types, each identified by a
numeric ID. Below is the complete list:

\bchdoc{btree-ids}

Btrees marked as ``snapshot-aware'' include a snapshot ID in every key position,
allowing different snapshot versions to coexist within the same btree. Some
btrees are ``write-buffered'': updates are batched in memory and flushed
periodically rather than applied directly to btree nodes, reducing write
amplification for high-frequency per-IO metadata (backpointers, LRU timestamps,
accounting counters). Write-buffered updates are unordered and
eventually-consistent: the btree does not reflect pending updates until the
next flush, and flushing sorts by key position, discarding temporal ordering.
Flushing deduplicates redundant updates, then walks btree nodes in order for
efficient bulk insertion; this sort-merge-sweep is inherently single-threaded,
making its efficiency critical for multithreaded workloads.

\subsection{Btree keys}

\bchdoc{bkey-structures}

See \S\ref{bkey-type-list} for the complete list of key types.

\subsection{Btree key types}
\label{bkey-type-list}

\bchdoc{bkey-types}

\subsection{Metadata versions}

Each on-disk format change is identified by a version number using
\texttt{BCH\_VERSION(major, minor)}. The version history records what changed
and when, serving as a changelog for the on-disk format. The filesystem's
superblock records both the current version and the minimum version of any data
still on disk, allowing compatibility code to be safely dropped.

\bchdoc{metadata-versions}

\end{document}
